\begin{proposition}\label{prop:bounded-energy-identities}
\Cref{def:model-energy} and
\eqref{eq:energy-cocycle}--\eqref{eq:energy-sandwich} extend to bounded
\(\omega\)-psh functions. They are continuous under decreasing bounded nets.
\end{proposition}

\begin{proof}
Let \(\varphi\) be a bounded \(\omega\)-psh function.  By
\Cref{thm:bfj-decreasing-regularization} choose a decreasing net
\(\varphi_\alpha\searrow\varphi\) of \(\omega\)-psh model functions.
Every member already satisfies \(\varphi_\alpha\geq\varphi\), so the
bounded limit supplies a common lower bound.  For an upper bound put
\(B=\sup_{X^{\An}}\varphi+1\).  The open sets
\[
 U_\alpha:=\{x\in X^{\An}:\varphi_\alpha(x)<B\}
\]
cover \(X^{\An}\): at each point the net decreases to a value strictly
less than \(B\).  They are open because model functions are continuous.
Compactness of \(X^{\An}\), by \Cref{thm:berk-proper-compact}, gives a
finite subcover \(U_{\alpha_1},\ldots,U_{\alpha_r}\).  Choose
\(\alpha_0\) above these finitely many indices.  Decreasingness then gives
\(\varphi_\alpha<B\) everywhere for every \(\alpha\geq\alpha_0\).
The tail \(\{\alpha:\alpha\geq\alpha_0\}\) is cofinal in the original
directed set, so restriction to it preserves the pointwise limit.  On
this tail there is a finite \(M\) with
\(-M\leq\varphi\leq\varphi_\alpha\leq M\).
Only indices have been discarded: no potential has been clipped or
translated, and all the approximants remain the given psh model functions.

\emph{Step 1: existence of the limit.}  Each summand of
\eqref{eq:model-energy} for \(\varphi_\alpha\) is of the form
\(\int h^\alpha\MA_\omega(\varphi_1^\alpha,\ldots,\varphi_n^\alpha)\) with
\(h^\alpha=\varphi_\alpha\searrow\varphi\) and every
\(\varphi_i^\alpha\in\{\varphi_\alpha,0\}\).  By
\eqref{eq:weighted-monotone-ma} each summand converges, so
\(E_\omega(\varphi_\alpha)\) converges; we define \(E_\omega(\varphi)\) to be
the limit and observe that it is given by \eqref{eq:model-energy}, the
integrals now being taken against the bounded mixed products of
\Cref{thm:bounded-ma-extension}.  In particular, the definition is finite
and agrees with the model energy when \(\varphi\) is a model function.

\emph{Step 2: independence of the net.}  Let \(\varphi_\beta'\searrow\varphi\)
be a second decreasing net of \(\omega\)-psh model functions, restricted
to a uniformly bounded cofinal tail as above.  Fix
\(\alpha\) and \(\varepsilon>0\).  The sets
\[
 F_\beta:=\{\varphi_\beta'\geq\varphi_\alpha+\varepsilon\}\subset X^{\An}
\]
are closed, because both functions are continuous, and decrease with
\(\beta\).  Their intersection is contained in
\(\{\varphi\geq\varphi_\alpha+\varepsilon\}=\varnothing\).
Compactness and directedness therefore make one of them empty, so
\(\varphi_\beta'\leq\varphi_\alpha+\varepsilon\) eventually in \(\beta\).
Monotonicity of \(E_\omega\) at model level
(\Cref{cor:energy-monotone-concave}) and
\eqref{eq:bfj-energy-translation} give
\(\lim_\beta E_\omega(\varphi_\beta')\leq
E_\omega(\varphi_\alpha)+\varepsilon V\).  Taking the limit in \(\alpha\)
and then \(\varepsilon\downarrow0\), and repeating with the two nets
exchanged, shows that the limits agree.  The translated function is used
only as a comparison majorant; neither approximating net is translated.

\emph{Step 3: the identities.}  Fix bounded \(\varphi,\psi\) and decreasing
model nets \(\varphi_\alpha\searrow\varphi\),
\(\psi_\beta\searrow\psi\).  Restrict each to a uniformly bounded cofinal
tail and use the product of the two directed sets.  Every integral in
\eqref{eq:energy-cocycle}--\eqref{eq:energy-sandwich} is of the form
\(\int h\,\MA_\omega(\chi_1,\ldots,\chi_n)\), with \(h\) a difference of
two members of decreasing bounded psh nets and each \(\chi_i\) a member
of such a net.  Splitting \(h=\psi_\beta-\varphi_\alpha\) into its two
terms and applying \eqref{eq:weighted-monotone-ma} to each, every such
integral converges to the corresponding integral for \(\varphi,\psi\).
Thus the model identities and inequalities pass to the limit.  For
\eqref{eq:energy-first-variation} and
\eqref{eq:energy-second-variation}, note that
\(\varphi_{t,\alpha\beta}=\varphi_\alpha+t(\psi_\beta-\varphi_\alpha)\)
decreases to \(\varphi_t\) for each fixed \(t\in[0,1]\).  The expansion
\eqref{eq:ch26-energy-taylor} expresses the energy as a polynomial in
\(t\) of degree at most \(n+1\).  Each coefficient is a finite linear
combination of the bounded mixed integrals just treated, so its limit
exists.  Differentiating these finite polynomial expansions therefore
commutes with passage to the limit and proves the two variation formulas.

\emph{Step 4: decreasing bounded nets.}  Let
\(\varphi_\gamma\searrow\varphi\) be a decreasing net of bounded
\(\omega\)-psh functions with bounded limit.  Fix one index \(\gamma_0\).
On its cofinal tail,
\[
 \inf_{X^{\An}}\varphi\leq\varphi\leq\varphi_\gamma
 \leq\varphi_{\gamma_0}\leq\sup_{X^{\An}}\varphi_{\gamma_0},
 \qquad \gamma\geq\gamma_0,
\]
so the entire tail is uniformly bounded.  By Step~1 the bounded energy
formula \eqref{eq:model-energy} already holds for each
\(\varphi_\gamma\) and for \(\varphi\).  Apply
\eqref{eq:weighted-monotone-ma} to each of its summands, with weight
\(h^\gamma=\varphi_\gamma\) and with every measure slot equal to either
\(\varphi_\gamma\) or \(0\).  The resulting limits give
\(E_\omega(\varphi_\gamma)\to E_\omega(\varphi)\).
The bounded energy inequalities from Step~3 imply monotonicity, hence in
fact \(E_\omega(\varphi_\gamma)\searrow E_\omega(\varphi)\).
This uses the bounded Monge--Amp\`ere continuity theorem already established
in \Cref{chap:bfj-ma}, not a new choice or diagonalization of model nets.
\end{proof}
